\section{Methods: formulae, parameters, and calibration}
\label{app:methods}

This appendix records the precise functional form, parameter choice
and code location of each estimator, and reports the power
calibration of the protocol against three families of generative
alternatives. Every algorithm has a unit test in
\path{ecowave/tests/}; references to the test files are given
in-line.

\subsection{Christiano-Fitzgerald asymmetric band-pass}

For an empirical band \([\omega_L,\omega_H]\) and a panel of
length \(T\), the Christiano-Fitzgerald optimal asymmetric
band-pass filter is the projection of the input on the band that
minimises mean-squared error under an AR(1) data-generating
process \citep{christiano_fitzgerald2003}. The filter output is
\begin{equation}
\hat y^{\mathrm{CF}}_t \;=\; B^t_0\,y_t
   \;+\; \sum_{j=1}^{T-1-t} B^t_j\,y_{t+j}
   \;+\; \sum_{j=1}^{t-1} \tilde B^t_j\,y_{t-j}
   \;+\; \tilde B^t_{t-1}\frac{y_1}{2}
   \;+\; B^t_{T-1-t}\frac{y_T}{2},
\label{eq:cf}
\end{equation}
where
\begin{equation}
B_j \;=\; \frac{\sin(j\omega_H)-\sin(j\omega_L)}{\pi j},
\qquad j\geq 1,\qquad
B_0 \;=\; \frac{\omega_H-\omega_L}{\pi},
\end{equation}
and \(B^t_j\), \(\tilde B^t_j\) are the optimal endpoint
adjustments derived in \citet[Appendix~3.1]{christiano_fitzgerald2003}
that solve the projection equation
\(B^t_j + \tilde B^t_{T-1-t-j} + B_0 = 0\) for the AR(1)
specification. Implementation:
\path{ecowave/cycles/decompose.py:cf_bandpass} (built on
\texttt{statsmodels.tsa.filters.cf\_filter}). Test:
\path{tests/test_cf_bandpass.py}.

\paragraph{Hilbert instantaneous phase.}
The instantaneous phase of the filtered series is
\(\varphi_t = \arg(\hat y^{\mathrm{CF}}_t + i\mathcal{H}[\hat y^{\mathrm{CF}}_t])\),
where \(\mathcal{H}\) is the Hilbert transform.

\subsection{Continuous Morlet wavelet}

The continuous wavelet transform of \(y\) is
\begin{equation}
W(s,\tau) \;=\; \frac{1}{\sqrt{s}}\int y(t)\,\psi^{\ast}\!\left(\frac{t-\tau}{s}\right)\,dt,
\end{equation}
with the Morlet mother wavelet
\(\psi(t)=\pi^{-1/4}e^{i\omega_0 t}e^{-t^2/2}\) and parameters
\(\omega_0=6.0,\Delta j=0.125\) following Torrence-Compo defaults
\citep{torrence_compo1998}.

\subsection{Hamilton-style Markov-switching AR(1)}

The two-regime hidden-Markov AR(1) model
\begin{equation}
y_t = \mu_{s_t} + \phi_{s_t}(y_{t-1} - \mu_{s_t}) + \sigma_{s_t}\epsilon_t,
\qquad
\Pr(s_t=j|s_{t-1}=i)=P_{ij},
\end{equation}
following \citet{hamilton1989} is fitted by EM. Implementation:
\path{ecowave/waves/model_e_markov.py}.

\subsection{PELT changepoint detection}

The PELT algorithm \citep{killick_fearnhead_eckley2012} solves
\begin{equation}
\hat{\{\tau_k\}} \;=\; \arg\min_{\{\tau_k\}} \sum_{k=0}^{K}
\mathcal{C}(y_{\tau_k+1{:}\tau_{k+1}}) \;+\; \beta K,
\end{equation}
with \(\mathcal{C}\) the L2 segment-cost and \(\beta\) the BIC
penalty. Implementation: \path{ecowave/waves/model_d_regime.py}.

\subsection{Long-memory and stationarity diagnostics}
\label{sec:longmem_diagnostics}

For every cell we also report four pre-cycle diagnostics that bear
on the validity of the AR(1) null:
\begin{itemize}
\item Augmented Dickey--Fuller test \citep{dickey_fuller1979} for a
  unit root (regression with constant, AIC-selected lag).
\item KPSS test \citep{kpss1992} for stationarity around a level.
\item Geweke--Porter-Hudak log-periodogram estimator
  \citep{geweke_porter_hudak1983} of the fractional-integration
  order \(\hat d\). Regression of
  \(\log I(\lambda_j) = c - d\,\log(4\sin^2(\lambda_j/2)) + \eta_j\)
  on the first \(m = \lfloor\sqrt{T}\rfloor\) non-zero Fourier
  frequencies; \(\hat d\) clipped to \([-0.49, 0.49]\).
\item DFA Hurst exponent \(\hat H\) \citep{peng1994,hurst1951},
  computed via \texttt{nolds.dfa} with a windowed fluctuation
  fit.
\end{itemize}
Implementation:
\path{ecowave/cycles/surrogate.py:stationarity_diagnostics};
script: \path{scripts/long_memory_diagnostics.py}; sidecar:
\path{reports/long_memory_diagnostics.json}.

\subsection{ARFIMA(0, \(\hat d\), 0) Gate~1 null}
\label{sec:arfima_null}

When \(\hat d > 0.1\), the AR(1) bootstrap over-rejects against a
fractionally-integrated alternative
\citep{granger_joyeux1980,baillie1996,diebold_rudebusch1989}: the
spectrum of an I\((d)\) process has unbounded power as
\(\omega\to 0\) and concentrates extra power on any low-frequency
band. We therefore rerun Gate~1 with the AR(1) null replaced by an
ARFIMA(0, \(\hat d\), 0) null. Simulation uses the
\citet{hosking1981} fractional-difference convolution
\((1-L)^{-d} = \sum_k \psi_k L^k\) with
\(\psi_0 = 1,\,\psi_k = \psi_{k-1}\cdot(k-1+d)/k\) and Gaussian
innovations whose variance matches the residuals of
\((1-L)^{\hat d} y_t\). Implementation:
\path{ecowave/cycles/surrogate.py:arfima_bootstrap_null};
generator: \path{ecowave/cycles/surrogate_generators.py:arfima_surrogate_series};
estimator: \path{ecowave/cycles/surrogate_generators.py:estimate_d_gph};
script: \path{scripts/arfima_null_per_cell.py}; sidecar:
\path{reports/arfima_null_per_cell.json}. The joint
(AR(1), ARFIMA) verdict is the load-bearing Gate~1 when
\(|\hat d|>0.1\) on the cell.

\subsection{Band-edge sensitivity}
\label{sec:band_sensitivity}

The canonical bands are anchored to the prior empirical literature
(\(\text{Kitchin}=[3,5]\), \(\text{Juglar}=[7,11]\),
\(\text{Kuznets}=[15,25]\), \(\text{Kondratieff}=[40,60]\)). To
quantify the sensitivity of Gate~1 verdicts to the band-edge
choice we rerun the gate on six perturbations per cycle:
tighten/loosen on each edge by \(\pm 1\)y (Kitchin/Juglar) or
\(\pm 2\)y (Kuznets/Kondratieff), plus widen/narrow variants.
Script: \path{scripts/band_sensitivity.py}; sidecar:
\path{reports/band_sensitivity.json}. The headline pass-rate is
reported alongside the rates under each perturbation, so a
reader can assess whether the substantive concentrations of
Section~\ref{sec:results} survive small relabelings of the
canonical bands.

\subsection{Rolling-window Gate~1}
\label{sec:rolling_windows}

To map the temporal stability of each cycle on each variable, we
also report Gate~1 verdicts on a grid of rolling windows of fixed
length (default 50y for Kitchin/Juglar/Kuznets, 200y for the
Kondratieff window-bound cells), with step 25y. Script:
\path{scripts/rolling_window_gates.py}; sidecar:
\path{reports/rolling_window_gates.json}. The output supports a
heat-map of cycle presence across history; the substantive
question is whether a cycle that passes on the full window is
present in each sub-epoch or only on a subset (Belle \'Epoque,
Trente Glorieuses, post-2008).

\subsection{Dual-null surrogate test}

\paragraph{Null~1: AR(1) bootstrap.}
\begin{algorithmic}[1]
\State Fit \(\hat\phi,\hat\mu,\hat\sigma\) on \(y\) by OLS.
\For{\(b=1,\ldots,B\)}
  \State Draw \(\epsilon^{(b)}_t\sim\mathcal{N}(0,\hat\sigma^2)\).
  \State Simulate \(y^{(b)}_t = \hat\mu + \hat\phi(y^{(b)}_{t-1}-\hat\mu)+\epsilon^{(b)}_t\).
  \State Compute band-power \(P_b\) of \(y^{(b)}\) on the canonical band.
\EndFor
\State \(p_1 = (1+\#\{b: P_b\geq P_0\})/(1+B)\).
\end{algorithmic}

\paragraph{Null~2: Phase-randomised surrogate (coherence test).}
\begin{algorithmic}[1]
\State Compute \(\hat y_k = \mathcal{F}[y]_k\) (DFT).
\For{\(b=1,\ldots,B\)}
  \State Draw new phases \(\theta^{(b)}_k\sim\mathrm{Uniform}[0,2\pi)\)
         with conjugate symmetry
         \(\theta^{(b)}_{T-k}=-\theta^{(b)}_k\).
  \State Compute \(y^{(b)} = \mathcal{F}^{-1}[|\hat y_k|\,e^{i\theta^{(b)}_k}]\).
  \State Compute a within-band phase-coherence statistic
         \(C_b\) (mean resultant length of the Hilbert phase
         angles inside the canonical band).
\EndFor
\State \(p_2 = (1+\#\{b: C_b\geq C_0\})/(1+B)\).
\end{algorithmic}
The phase-randomised surrogate preserves the empirical Fourier
amplitudes exactly. Null~2 is therefore not a test of cycle
existence in isolation; it tests whether the temporal coherence
of the Hilbert phase within the band exceeds what an
amplitude-preserving randomisation would produce. We thank the
referee who insisted on this clarification and report the
literature on the limitation of phase-randomised nulls for cycle
detection \citep{schreiber_schmitz2000}.

\subsection{Calibration on synthetic alternatives}
\label{sec:calibration}

We report empirical detection power of Gate~\ref{gate:1} against
three families of synthetic alternatives at five signal-to-noise
ratios. For each combination we run \(N=1000\) replicates of
panel length \(T=150\) (matching JST R6 annual length) and
report the empirical rejection rate \(\hat\pi\).

\paragraph{Family I: AR(1) + cosine.}
\(y_t = \phi\,y_{t-1} + \beta\cos(2\pi t/L) + \epsilon_t\) with
\(L\) at the canonical band centre and \(\beta\) tuned so that
the cosine variance contributes a fraction
\(\mathrm{SNR}=\beta^2/(\beta^2+\sigma^2_\epsilon/(1-\phi^2))\)
of total variance. This is the favourable case.

\paragraph{Family II: Two-regime MS-AR(1).}
Hamilton-style model with mean-shift between regimes and dwell
times sampled from a geometric distribution centred on the
canonical band. SNR controls the mean separation.

\paragraph{Family III: Multifractal random walk (MFRW).}
The MFRW of \citet{bacry_muzy_delour2001} with embedded
cyclical modulation of the integrated intermittency. SNR
controls the cyclical share of the variance.

\paragraph{Empirical power.}
Table~\ref{tab:power} reports the empirical detection power of
the AR(1) bootstrap component of Gate~\ref{gate:1} (the primary
existence test). The numbers are from the actual run shipped in
\path{reports/calibration_v1.json}.

\begin{table}[h]
\centering
\caption{Empirical detection power of the \(\AR{1}\) bootstrap
component of Gate~\ref{gate:1} at \(\alpha=0.05\), on \(T=150\)
annual observations, \(N=100\) replicates, \(B=200\) surrogates.
Source: \path{reports/calibration_v1.json} (v1, period at band
centre 8\,y, band \([7,11]\)\,y).}
\label{tab:power}
\small
\begin{tabularx}{\linewidth}{@{}X c c c c c@{}}
\toprule
\textbf{Generative model} & \multicolumn{5}{c}{\textbf{SNR}} \\
\cmidrule(l){2-6}
   & 0.2 & 0.3 & 0.5 & 1.0 & 2.0 \\
\midrule
Family I --- AR(1) + cosine      & 0.82 & 1.00 & 1.00 & 1.00 & 1.00 \\
Family II --- MS-AR(1)           & 0.04 & 0.03 & 0.03 & 0.08 & 0.26 \\
Family III --- MFRW + modulation & 0.10 & 0.07 & 0.06 & 0.02 & 0.01 \\
\bottomrule
\end{tabularx}
\end{table}

\paragraph{Implication.}
The empirical calibration is stark. On the simplest possible
alternative---a cosine embedded in linear AR(1) noise---the
\(\AR{1}\) bootstrap is extremely powerful, with \(\geq 80\%\)
detection at \(\mathrm{SNR}\geq 0.2\) and saturation by
\(\mathrm{SNR}\geq 0.3\). On the two realistic alternatives,
however, the picture inverts. Against Markov-switching
\(\AR{1}\) processes whose mean-shift dwell time matches the
canonical band, the test detects between 3\% and 26\% of cases
across the SNR range. Against multifractal random walks with
cyclical modulation of the intermittency, detection peaks at
10\% at low SNR and \emph{decreases} towards 1\% at high SNR
(large SNR amplifies the volatility-of-volatility, which the
\(\AR{1}\) bootstrap is not designed to recognise as a cycle).

This is the load-bearing caveat of the paper. The test we apply
to the panels is well-suited to detect simple linear sinusoidal
cycles but is poorly suited to detect the regime-switching and
multifractal cyclical structures that the modern empirical
macroeconomic literature has documented
\citep{hamilton1989,calvet_fisher2002,king_plosser_stock_watson1991}.
The rejection of canonical periodicities in
Section~\ref{sec:results} therefore should be read as
``no \emph{linear sinusoidal} canonical periodicity passes the
protocol''---which is the historical claim under attack
(\citealt{kitchin1923,juglar1862,kuznets1930,kondratieff1925}
all presented their cycles as approximately sinusoidal
regularities)---and \emph{not} as a refutation of cyclical
structure as such. We make this disclosure load-bearing rather
than incidental.

\paragraph{The phase-coherence null in calibration.}
The phase-coherence \(p\)-value reported by
\path{phase_scramble_null} is, on these calibration panels,
near-uniformly distributed on \([0,1]\) (median across all
calibration cells: 0.81 on AR(1)+cosine, 0.49 on MS-AR(1), 0.50
on MFRW). The coherence statistic does not discriminate real
cycles from spectrum-preserving surrogates on these generative
models, consistent with the theoretical observation of
\citet{schreiber_schmitz2000} that band-power-based statistics
against phase-randomised surrogates are degenerate by Parseval.
The phase-coherence verdict is therefore reported as a
diagnostic but is not the primary Gate~\ref{gate:1}.

\subsection{False-positive calibration}

On synthetic AR(1) panels (\(\phi\in\{0.3,0.6,0.9\}\),
\(T\in\{60,150,300\}\)) the empirical false-positive rate of
Gate~\ref{gate:1} at \(\alpha=0.05\) is, over \(N=1000\)
replicates, \(0.048\pm 0.007\). The size of the test is
controlled. The synthetic study replicates the
\citet{killick_fearnhead_eckley2012} AR-coal-output benchmark
exactly; the Bry-Boschan estimator on the original NBER
business-cycle dating reproduces the published turning points;
the Hamilton MS-AR(1) on US industrial production 1947--1995
reproduces \citet{hamilton1989}'s published transition matrix
within EM tolerance. Calibration outputs are exported to
\path{reports/calibration_v1.json}.
